%%
%% Concrete Implementation Proposal: Numeric Tower Phase 1
%%
%% Companion to "Design Considerations for a Full Numeric Tower in MScheme"
%%
\documentclass[12pt]{article}
\usepackage[margin=1in]{geometry}
\usepackage{booktabs}
\usepackage{listings}
\usepackage{xcolor}
\usepackage{hyperref}
\usepackage{textcomp}
\usepackage{amssymb}
\usepackage{amsmath}
\usepackage{enumitem}

\lstdefinelanguage{Scheme}{
  morekeywords={define,lambda,if,cond,else,let,let*,letrec,begin,
    set!,quote,quasiquote,unquote,unquote-splicing,and,or,not,
    case,do,delay,force,map,for-each,apply,call-with-current-continuation,
    call/cc,require-modules,load,macro,define-global-symbol,
    exact->inexact,inexact->exact,make-rectangular,make-polar,
    real-part,imag-part,magnitude,angle,numerator,denominator,
    rationalize,number->string,string->number},
  sensitive=true,
  morecomment=[l]{;},
  morestring=[b]",
  literate={lambda}{$\lambda$}1,
}

\lstdefinelanguage{Modula3}{
  morekeywords={MODULE,INTERFACE,IMPORT,FROM,PROCEDURE,BEGIN,END,
    VAR,CONST,TYPE,RECORD,OBJECT,METHODS,OVERRIDES,RETURN,
    IF,THEN,ELSE,ELSIF,WHILE,DO,FOR,TO,BY,LOOP,EXIT,REPEAT,UNTIL,
    CASE,OF,WITH,TRY,EXCEPT,FINALLY,RAISE,RAISES,
    NEW,ARRAY,REF,REFANY,SET,BOOLEAN,INTEGER,REAL,LONGREAL,EXTENDED,
    CHAR,TEXT,LONGINT,
    TRUE,FALSE,NIL,BRANDED,REVEAL,EVAL,NOT,AND,OR,IN,DIV,MOD,
    ISTYPE,NARROW,TYPECODE,GENERIC,ADDRESS,BITS,BITSIZE},
  sensitive=true,
  morecomment=[n]{(*}{*)},
  morestring=[b]",
}

\lstset{
  basicstyle=\ttfamily\small,
  breaklines=true,
  frame=single,
  backgroundcolor=\color{gray!5},
  keywordstyle=\bfseries\color{blue!70!black},
  commentstyle=\itshape\color{green!50!black},
  stringstyle=\color{red!60!black},
  showstringspaces=false,
  columns=flexible,
  tabsize=2,
  xleftmargin=2em,
  framexleftmargin=1.5em,
}

\title{\bfseries\sffamily\fontsize{24.88}{30}\selectfont
  Numeric Tower Phase~1:\\Implementation Proposal}
\author{Mika Nystr\"om \\ {\tt mika@alum.mit.edu}}
\date{March 2026}

\begin{document}
\maketitle
\thispagestyle{empty}

\newpage
\tableofcontents
\newpage

%======================================================================
\section{Overview}
\label{sec:overview}
%======================================================================

This document is the companion to \emph{Design Considerations for a
Full Numeric Tower in MScheme}, which analyzes the design space.
Here we give a concrete, code-level implementation plan for Phase~1:
exact integers (fixnum and bignum).  Sections are organized by
subsystem, each listing the files to create or modify, the
interfaces involved, and Modula-3 code sketches.  The overall goal
is R5RS-conformant exact/inexact integer semantics with minimal
disruption to the existing codebase.

Phase~1 is self-contained and does not require Phases~2--4
(rationals, complex numbers, compiler unboxing optimizations).

%======================================================================
\section{Build Instructions}
\label{sec:build}
%======================================================================

All MScheme packages live under {\tt m3-scheme/} in the CM3 tree.
The CM3 compiler must be on your path (or referenced by absolute
path).

%----------------------------------------------------------------------
\subsection{Building MScheme}
%----------------------------------------------------------------------

\begin{lstlisting}[language={}]
# Set up
export CM3=/Users/mika/cm3/install/bin/cm3

# Build the mscheme library
cd m3-scheme/mscheme
$CM3 -build
$CM3 -ship

# Build the mscheme-interactive binary
cd ../mscheme-interactive
$CM3 -build

# The binary is at:
#   mscheme-interactive/ARM64_DARWIN/mscheme
\end{lstlisting}

\noindent \textbf{Important:} after modifying {\tt mscheme} (the
library), you must {\tt \$CM3 -ship} it before rebuilding
{\tt mscheme-interactive}, because the binary links against the
shipped shared library ({\tt libmscheme.5.dylib}), not the local
{\tt .a} archive.  If the binary appears not to pick up your
changes, this is almost certainly the reason.

%----------------------------------------------------------------------
\subsection{Running Tests}
%----------------------------------------------------------------------

\begin{lstlisting}[language={}]
cd m3-scheme/mscheme-doc

# Numeric tower tests (Phase 1 target: 158/158)
mscheme test-numeric-tower.scm

# R4RS conformance tests (must stay at 95/95)
mscheme test-r4rs-gaps.scm

# Existing manual examples
mscheme test-examples.scm
\end{lstlisting}

%----------------------------------------------------------------------
\subsection{Building the Compiler}
%----------------------------------------------------------------------

The MScheme compiler is a separate package that depends on the
mscheme library.

\begin{lstlisting}[language={}]
# Build the compiler package
cd m3-scheme/schemecompiler
$CM3 -build
$CM3 -ship

# Build the compiled-mscheme binary
cd ../mscheme-interactive_r
$CM3 -build
\end{lstlisting}

\noindent The compiled binary is at
{\tt mscheme-interactive\_r/ARM64\_DARWIN/mscheme}.

%----------------------------------------------------------------------
\subsection{Bootstrap Rebuild}
%----------------------------------------------------------------------

When modifying the compiler itself ({\tt compiler.scm}), a
two-phase bootstrap is needed:

\begin{lstlisting}[language={}]
# Phase B: compile new compiler under old compiler
cd m3-scheme/schemecompiler
$CM3 -build && $CM3 -ship
cd ../mscheme-interactive_r
$CM3 -build

# Phase C: recompile compiler under itself
cd ../schemecompiler
$CM3 -build && $CM3 -ship
cd ../mscheme-interactive_r
$CM3 -build
\end{lstlisting}

%======================================================================
\section{New {\tt SchemeInt} Interface and Module}
\label{sec:impl-schemeint}
%======================================================================

A new module {\tt SchemeInt} encapsulates all exact-integer
operations.  It is the single point of contact for creating,
testing, and operating on fixnums and bignums.

\paragraph{Files.}
\begin{itemize}[nosep]
\item {\tt mscheme/src/SchemeInt.i3} (new)
\item {\tt mscheme/src/SchemeInt.m3} (new)
\end{itemize}

\medskip\noindent\textbf{Interface.}
\begin{lstlisting}[language=Modula3]
INTERFACE SchemeInt;
IMPORT SchemeObject;

TYPE T = REF INTEGER;  (* fixnum *)

PROCEDURE FromI(x: INTEGER): SchemeObject.T;
  (* Cache [-256..256]; allocate outside that range. *)

PROCEDURE IsInt(x: SchemeObject.T): BOOLEAN;
  (* TRUE iff x is REF INTEGER. *)

PROCEDURE IsExactInt(x: SchemeObject.T): BOOLEAN;
  (* TRUE iff x is REF INTEGER or Mpz.T. *)

PROCEDURE ToInteger(x: SchemeObject.T): INTEGER;
  (* Extract value from REF INTEGER.  Checked runtime
     error if x is not REF INTEGER. *)

PROCEDURE Add(a, b: INTEGER): SchemeObject.T;
PROCEDURE Sub(a, b: INTEGER): SchemeObject.T;
PROCEDURE Mul(a, b: INTEGER): SchemeObject.T;
  (* Overflow-checked via Word; promote to Mpz.T
     on overflow. *)

END SchemeInt.
\end{lstlisting}

\paragraph{Overflow detection.}
CM3 integer overflow is undefined behavior (see the design report,
Section~5.2).  The fixnum arithmetic procedures must therefore
detect overflow \emph{before} it occurs, or use {\tt Word}-based
unsigned arithmetic with post-hoc sign reconstruction.  A robust
approach uses the {\tt Word} module:

\begin{lstlisting}[language=Modula3]
PROCEDURE Add(a, b: INTEGER): SchemeObject.T =
  VAR sum := Word.Plus(a, b);
      sa  := Word.And(a, SignBit);
      sb  := Word.And(b, SignBit);
      ss  := Word.And(sum, SignBit);
  BEGIN
    (* Overflow iff both operands have the same sign
       and the result has a different sign. *)
    IF Word.And(Word.Not(Word.Xor(sa, sb)),
                Word.Xor(sa, ss)) # 0 THEN
      (* Promote to bignum *)
      WITH ma = Mpz.InitSet(a), mb = Mpz.InitSet(b) DO
        Mpz.Add(ma, ma, mb);
        RETURN ma
      END
    ELSE
      RETURN FromI(LOOPHOLE(sum, INTEGER))
    END
  END Add;
\end{lstlisting}

\noindent Subtraction is analogous (overflow when signs differ and
result sign disagrees with the first operand).  Multiplication
requires a wider intermediate result; on 64-bit targets the
simplest correct approach is to promote both operands to {\tt
Mpz.T}, multiply, and demote if the result fits:

\begin{lstlisting}[language=Modula3]
PROCEDURE Mul(a, b: INTEGER): SchemeObject.T =
  VAR ma := Mpz.InitSet(a);
      mb := Mpz.InitSet(b);
  BEGIN
    Mpz.Mul(ma, ma, mb);
    IF Mpz.FitsInteger(ma) THEN
      RETURN FromI(Mpz.ToInteger(ma))
    ELSE
      RETURN ma
    END
  END Mul;
\end{lstlisting}

\noindent This is not the fastest possible approach---for
multiplication, a hardware-level widening multiply would be
preferable---but it is correct and simple.  Optimization can follow
once the tower is functional.

\paragraph{Fixnum cache.}
The cache is initialized at module load time:

\begin{lstlisting}[language=Modula3]
VAR cache: ARRAY [-256..256] OF SchemeObject.T;

PROCEDURE FromI(x: INTEGER): SchemeObject.T =
  BEGIN
    IF x >= -256 AND x <= 256 THEN
      RETURN cache[x]
    ELSE
      WITH r = NEW(REF INTEGER) DO r^ := x; RETURN r END
    END
  END FromI;

BEGIN
  FOR i := -256 TO 256 DO
    WITH r = NEW(REF INTEGER) DO
      r^ := i; cache[i] := r
    END
  END
END SchemeInt.
\end{lstlisting}

%======================================================================
\section{Modified {\tt SchemeLongReal} (Backward-Compatible Bridge)}
\label{sec:impl-longreal}
%======================================================================

The existing {\tt SchemeLongReal.FromO} procedure is the
backward-compatibility linchpin: all existing Modula-3 code that
consumes Scheme numbers calls it.  Two procedures need extension.

\paragraph{Files.}
\begin{itemize}[nosep]
\item {\tt mscheme/src/SchemeLongReal.m3} (modify)
\end{itemize}

\medskip\noindent\textbf{{\tt FromO} extension.}
\begin{lstlisting}[language=Modula3]
PROCEDURE FromO(x: SchemeObject.T): LONGREAL =
  BEGIN
    TYPECASE x OF
    | REF INTEGER (ri) => RETURN FLOAT(ri^, LONGREAL);
    | Mpz.T (m)        => RETURN Mpz.GetD(m);
    | SchemeLongReal.T (lr) => RETURN lr.val;
    ELSE
      RAISE E("not a number")
    END
  END FromO;
\end{lstlisting}

\noindent The two new branches ({\tt REF INTEGER} and {\tt Mpz.T})
are inserted \emph{before} the existing {\tt SchemeLongReal.T}
branch.  Existing callers that pass {\tt SchemeLongReal.T} values
are unaffected.

\medskip\noindent\textbf{{\tt Int} extension.}
\begin{lstlisting}[language=Modula3]
PROCEDURE Int(x: SchemeObject.T): INTEGER =
  BEGIN
    TYPECASE x OF
    | REF INTEGER (ri) => RETURN ri^;
    | Mpz.T (m)        => RETURN Mpz.ToInteger(m);
    | SchemeLongReal.T (lr) =>
        RETURN ROUND(lr.val)
    ELSE
      RAISE E("not a number")
    END
  END Int;
\end{lstlisting}

%======================================================================
\section{Reader Changes ({\tt SchemeInputPort.m3})}
\label{sec:impl-reader}
%======================================================================

The reader currently parses all numeric literals as {\tt LONGREAL}.
After the change, the parser distinguishes integer literals from
floating-point literals.

\paragraph{Files.}
\begin{itemize}[nosep]
\item {\tt mscheme/src/SchemeInputPort.m3} (modify)
\end{itemize}

\paragraph{Dispatch rule.}
\begin{enumerate}[nosep]
\item If the token contains a decimal point or exponent marker
  ({\tt e}, {\tt E}, {\tt s}, {\tt f}, {\tt d}, {\tt l}),
  parse as {\tt LONGREAL} via {\tt Scan.LongReal} and return
  {\tt SchemeLongReal.FromLR(...)}.  This path is unchanged.

\item If the token is a pure integer (digits with optional sign,
  no~`.', no exponent), parse as {\tt INTEGER} via {\tt Scan.Int}
  and return {\tt SchemeInt.FromI(...)}.  If {\tt Scan.Int} raises
  a range error (value exceeds {\tt LAST(INTEGER)}), fall through
  to the bignum path: parse with {\tt Mpz.InitSetStr}.

\item The {\tt \#e} (exact) prefix forces the integer parse path
  even when the token would otherwise parse as a float:
  {\tt \#e3.14} $\to$ {\tt SchemeInt.FromI(3)} (truncation).

\item The {\tt \#i} (inexact) prefix forces the {\tt LONGREAL} path
  even for integer tokens: {\tt \#i42} $\to$
  {\tt SchemeLongReal.FromLR(42.0d0)}.

\item Non-decimal bases ({\tt \#b}, {\tt \#o}, {\tt \#x}) continue
  to use {\tt Scan.Int} with the appropriate base parameter, now
  returning {\tt SchemeInt.FromI}.

\item R7RS special float literals: {\tt +inf.0}, {\tt -inf.0},
  {\tt +nan.0} are recognized and return the corresponding
  IEEE~754 values.  Also accept {\tt Infinity}, {\tt -Infinity},
  {\tt NaN} (CM3 {\tt Fmt} convention) for backward compatibility.
\end{enumerate}

%======================================================================
\section{Printer Changes ({\tt SchemeUtils.m3})}
\label{sec:impl-printer}
%======================================================================

The {\tt StringifyB} procedure (and its callers {\tt Stringify},
{\tt Display}) must recognize the new types.

\paragraph{Files.}
\begin{itemize}[nosep]
\item {\tt mscheme/src/SchemeUtils.m3} (modify)
\end{itemize}

\paragraph{New {\tt TYPECASE} branches.}
Insert before the existing {\tt SchemeLongReal.T} branch:

\begin{lstlisting}[language=Modula3]
| REF INTEGER (ri) =>
    Wr.PutText(wr, Fmt.Int(ri^))

| Mpz.T (m) =>
    Wr.PutText(wr, Mpz.Format(m, base := 10))
\end{lstlisting}

\noindent Exact integers print without a decimal point, satisfying
the R5RS requirement that {\tt (display 42)} prints {\tt 42}, not
{\tt 42.0} or {\tt 42.}.

%======================================================================
\section{Predicate Changes ({\tt SchemePrimitive.m3})}
\label{sec:impl-predicates}
%======================================================================

R5RS defines a hierarchy of numeric predicates with precise
semantics.  The current implementation collapses all of them to
``is it a {\tt SchemeLongReal.T}?''  Table~\ref{tab:predicates}
shows the required dispatch.

\begin{table}[htbp]
\centering
\begin{tabular}{lccc}
\toprule
\textbf{Predicate} & {\tt REF INTEGER} & {\tt Mpz.T} &
  {\tt SchemeLongReal.T} \\
\midrule
{\tt number?}   & \checkmark & \checkmark & \checkmark \\
{\tt complex?}  & \checkmark & \checkmark & \checkmark \\
{\tt real?}     & \checkmark & \checkmark & \checkmark \\
{\tt rational?} & \checkmark & \checkmark & if integer-valued \\
{\tt integer?}  & \checkmark & \checkmark & if integer-valued \\
{\tt exact?}    & \checkmark & \checkmark &  \\
{\tt inexact?}  &            &            & \checkmark \\
\bottomrule
\end{tabular}
\caption{Predicate dispatch for the three-level numeric tower.
  ``If integer-valued'' tests {\tt x = FLOOR(x)} and finiteness.}
\label{tab:predicates}
\end{table}

\paragraph{Files.}
\begin{itemize}[nosep]
\item {\tt mscheme/src/SchemePrimitive.m3} (modify)
\end{itemize}

%======================================================================
\section{Arithmetic Dispatch ({\tt SchemePrimitive.m3})}
\label{sec:impl-arith}
%======================================================================

The central arithmetic dispatcher, currently {\tt NumCompute},
must be rewritten to handle the three-level tower.

\paragraph{Dispatch logic for binary operations.}
\begin{enumerate}[nosep]
\item If both operands are {\tt REF INTEGER}: extract values,
  call {\tt SchemeInt.Add}/{\tt Sub}/{\tt Mul} (with overflow
  promotion).

\item If both operands are exact integers but at least one is
  {\tt Mpz.T}: promote the fixnum to {\tt Mpz.T}, perform the
  GMP operation, demote to fixnum if the result fits.

\item If any operand is {\tt SchemeLongReal.T} (inexact): coerce
  all operands to {\tt LONGREAL} via {\tt SchemeLongReal.FromO},
  perform the operation, return {\tt SchemeLongReal.FromLR(result)}.
  This is the \emph{exactness contagion} rule.
\end{enumerate}

\paragraph{Division.}
The {\tt /} operator has special semantics:
\begin{itemize}[nosep]
\item If both operands are exact integers and the division is
  exact (zero remainder), return an exact integer.
\item If both operands are exact integers and the division is
  not exact, return an inexact {\tt LONGREAL}.  (Phase~2 would
  return an exact rational instead.)
\item If any operand is inexact, return inexact.
\end{itemize}

\paragraph{Integer-only operations.}
{\tt quotient}, {\tt remainder}, and {\tt modulo} require exact
integer operands and signal an error on inexact arguments.  {\tt
gcd} and {\tt lcm} similarly operate on exact integers, using
{\tt Mpz.Gcd} and {\tt Mpz.Lcm} for the bignum case.

\paragraph{Comparison dispatch.}
{\tt NumCompare} follows the same three-level dispatch.  For
mixed exact/inexact comparisons, the exact operand is converted
to {\tt LONGREAL} before comparison.  (A more precise approach
would compare without conversion, but this suffices for Phase~1.)

%======================================================================
\section{{\tt eqv?}\ Changes ({\tt SchemeUtils.m3})}
\label{sec:impl-eqv}
%======================================================================

R5RS specifies that {\tt eqv?}\ on numbers compares both value
\emph{and} exactness.  In particular:

\begin{lstlisting}[language=Scheme]
(eqv? 1 1)     ;=> #t
(eqv? 1 1.0)   ;=> #f  (different exactness)
(eqv? 1.0 1.0) ;=> #t
\end{lstlisting}

\paragraph{Files.}
\begin{itemize}[nosep]
\item {\tt mscheme/src/SchemeUtils.m3} (modify)
\end{itemize}

\paragraph{New {\tt TYPECASE} branch.}
In the {\tt EqvP} procedure (or equivalent), add:
\begin{lstlisting}[language=Modula3]
IF ISTYPE(a, REF INTEGER) AND ISTYPE(b, REF INTEGER) THEN
  RETURN NARROW(a, REF INTEGER)^ = NARROW(b, REF INTEGER)^
ELSIF SchemeInt.IsExactInt(a) AND SchemeInt.IsExactInt(b) THEN
  (* Both exact; compare via Mpz if needed *)
  RETURN MpzEqual(a, b)
ELSIF SchemeInt.IsExactInt(a) # SchemeInt.IsExactInt(b) THEN
  (* One exact, one inexact: always #f *)
  RETURN FALSE
END
\end{lstlisting}

%======================================================================
\section{Compiler Changes ({\tt compiler.scm})}
\label{sec:impl-compiler}
%======================================================================

The MScheme compiler emits Modula-3 code.  Its numeric code
generation must be updated to use {\tt SchemeInt} for exact
integer constants and to add fixnum fast paths for inline
arithmetic.

\paragraph{Files.}
\begin{itemize}[nosep]
\item {\tt schemecompiler/compiler.scm} (modify)
\end{itemize}

\subsection{Constant References}

The {\tt constant-ref} procedure currently emits
{\tt SchemeLongReal.FromI($n$)} for all integer constants.
After the change:

\begin{itemize}[nosep]
\item Integer constants $\to$ {\tt SchemeInt.FromI($n$)}
\item Floating-point constants $\to$ {\tt SchemeLongReal.FromLR($n$d0)}
  (unchanged)
\end{itemize}

\noindent The module's {\tt IMPORT} list must be extended to include
{\tt SchemeInt}.

\subsection{Inline Arithmetic Fast Path}

The compiler currently inlines arithmetic operations by unboxing
to {\tt LONGREAL}.  The fixnum fast path wraps the existing code:

\begin{lstlisting}[language=Modula3]
IF ISTYPE(a, SchemeInt.T) AND ISTYPE(b, SchemeInt.T) THEN
  result := SchemeInt.Add(
    NARROW(a, SchemeInt.T)^,
    NARROW(b, SchemeInt.T)^)
ELSE
  result := SchemePrimitive.NumAdd(a, b)
END
\end{lstlisting}

\paragraph{Gradual overflow.}
The {\tt ELSE} branch must \emph{not} unconditionally coerce to
{\tt LONGREAL}.  Consider an accumulator loop that gradually
overflows fixnum range: {\tt SchemeInt.Add} promotes the result
to {\tt Mpz.T}, but on the next iteration that bignum reaches
the {\tt ELSE} branch.  If the {\tt ELSE} branch converts to
{\tt LONGREAL} via {\tt SchemeLongReal.FromO}, exactness is
silently lost.  Instead, the {\tt ELSE} branch must call the
interpreter's full arithmetic dispatch (here shown as
{\tt SchemePrimitive.NumAdd}), which handles all three levels
of the tower correctly: fixnum~$\times$~bignum,
bignum~$\times$~bignum, and any~$\times$~inexact.

This means the inline fast path is a \emph{fixnum-only} fast
path.  When either operand is not a fixnum, the general dispatch
handles it.  This is slightly slower than the old all-{\tt LONGREAL}
path for mixed-type cases, but it is \emph{correct}---exactness
is never silently lost.

\paragraph{Self-tail-call loops.}
The existing {\tt LONGREAL} unboxing in self-tail-call
loops is \textbf{not changed} in Phase~1.  This means compiled
tight loops that the compiler has already identified as
{\tt LONGREAL}-typed still operate at native {\tt LONGREAL} speed.
The fixnum fast path applies only to the general-case inline
arithmetic (non-loop contexts, or loops where the compiler
cannot prove {\tt LONGREAL} typing).

\subsection{Inline Predicates}

Type predicates can be inlined more precisely:
\begin{itemize}[nosep]
\item {\tt (number?\ $x$)} $\to$ {\tt ISTYPE($x$, SchemeInt.T)} or
  {\tt ISTYPE($x$, Mpz.T)} or\\
  {\tt ISTYPE($x$, SchemeLongReal.T)}
\item {\tt (integer?\ $x$)} $\to$ {\tt ISTYPE($x$, SchemeInt.T)} or
  {\tt ISTYPE($x$, Mpz.T)}
\item {\tt (exact?\ $x$)} $\to$ same as {\tt integer?} (Phase~1)
\end{itemize}

%======================================================================
\section{Build System ({\tt m3makefile})}
\label{sec:impl-buildsys}
%======================================================================

\paragraph{Files.}
\begin{itemize}[nosep]
\item {\tt mscheme/src/m3makefile} (modify)
\end{itemize}

\medskip\noindent\textbf{Changes.}
\begin{lstlisting}
Module("SchemeInt")
import("mpz")      % for Mpz.T access
\end{lstlisting}

\noindent The {\tt mpz} package is already available in the CM3
tree ({\tt m3-libs/mpz}).  The {\tt import} directive makes its
interface visible to the mscheme package.

%======================================================================
\section{Bootstrap Sequence}
\label{sec:impl-bootstrap}
%======================================================================

Because the MScheme compiler is self-hosting (written in Scheme,
compiled by itself), introducing a new numeric representation
requires a careful bootstrap.  The sequence has three phases:

\paragraph{Phase~A: Interpreter only.}
Add the {\tt SchemeInt} module and modify the interpreter's reader,
printer, predicates, and arithmetic dispatch.  The compiler
continues to emit {\tt SchemeLongReal}-based code.  Test the
interpreter thoroughly: all existing tests must pass (with updated
expected outputs for integer printing), and new exact-integer tests
must pass.

\paragraph{Phase~B: Compiler modification.}
Modify {\tt compiler.scm} to emit {\tt SchemeInt}-based code (as
described in Section~\ref{sec:impl-compiler}).  Compile the new
compiler using the \emph{old} compiler (which still uses
{\tt LONGREAL} semantics internally).  The resulting binary is a
``Phase~B compiler'' that emits numeric-tower code but was itself
compiled under {\tt LONGREAL} semantics.  Test this compiler.

\paragraph{Phase~C: Self-compilation.}
Use the Phase~B compiler to compile itself.  The resulting binary
is now a fully numeric-tower-aware compiler, compiled under
numeric-tower semantics.  This is the production compiler.

\medskip\noindent
The two-phase bootstrap (B~$\to$~C) is necessary because the
compiler's own intermediate computations (e.g., counting
arguments, computing offsets) will behave differently under exact
integer semantics than under {\tt LONGREAL} semantics---even
though the results should be identical for values within the
{\tt LONGREAL} integer range.  The Phase~C recompilation confirms
that the compiler is a fixed point.

%======================================================================
\section{What Phase~1 Does \emph{Not} Change}
\label{sec:impl-not}
%======================================================================

\begin{itemize}
\item \textbf{No rationals.}  Exact division of non-divisible
  integers returns inexact {\tt LONGREAL}.  Phase~2 would
  introduce {\tt BigRational} or {\tt Mpq.T}.

\item \textbf{No complex numbers.}  Phase~3 would add
  {\tt SchemeComplex.T}.

\item \textbf{No compiler {\tt INTEGER} unboxing.}  The compiler's
  self-tail-call optimization continues to unbox to {\tt LONGREAL}.
  Adding {\tt INTEGER} unboxing is a Phase~4 optimization.

\item \textbf{No SRFI-151 bitwise operations.}  However, the
  foundation is laid: {\tt Mpz.T} supports two's complement
  bitwise operations natively.  Adding {\tt bitwise-and},
  {\tt bitwise-or}, {\tt arithmetic-shift}, etc.\ is
  straightforward once exact integers exist.
\end{itemize}

%======================================================================
\section{Ancillary Fixes}
\label{sec:impl-ancillary}
%======================================================================

Several pre-existing conformance issues surfaced during test
development.  These are independent of the exact-integer work but
should be fixed as part of Phase~1.

\subsection{Infinity and NaN Read Syntax (R7RS)}

MScheme prints IEEE~754 special values using Modula-3's {\tt Fmt}
conventions: {\tt Infinity}, {\tt -Infinity}, {\tt NaN}.  However,
the reader cannot parse these back---they are treated as symbols.
R7RS (\S7.1.1) specifies the syntax {\tt +inf.0}, {\tt -inf.0},
{\tt +nan.0}.

\paragraph{Changes.}
\begin{itemize}[nosep]
\item \textbf{Reader} ({\tt SchemeInputPort.m3}): recognize
  {\tt +inf.0}, {\tt -inf.0}, {\tt +nan.0} as numeric literals
  returning the corresponding IEEE~754 {\tt LONGREAL} values.
  Also accept {\tt Infinity}, {\tt -Infinity}, {\tt NaN} for
  backward compatibility with existing output.
\item \textbf{Printer} ({\tt SchemeUtils.m3}): change the
  output of special float values from {\tt Infinity}/{\tt NaN}
  to {\tt +inf.0}/{\tt -inf.0}/{\tt +nan.0} (R7RS style).
  This ensures that {\tt (read (write x))} round-trips.
\end{itemize}

\noindent No {\tt .scm} file in the CM3 tree currently references
{\tt Infinity}, {\tt NaN}, or the R7RS forms, so the output change
is safe.

\subsection{Banker's Rounding ({\tt round})}

R5RS~\S6.2.5 specifies that {\tt round} uses half-to-even
(banker's rounding): {\tt (round 2.5)} $\to$ {\tt 2.0}.
MScheme delegates to Modula-3's {\tt ROUND}, which uses
half-away-from-zero: {\tt (round 2.5)} $\to$ {\tt 3.0}.

\paragraph{Fix.}
Replace the {\tt P.Round} implementation in {\tt SchemePrimitive.m3}
with explicit half-to-even logic:

\begin{lstlisting}[language=Modula3]
P.Round =>
  WITH d = FromO(x),
       r = FLOAT(ROUND(d), LONGREAL) DO
    (* If exactly halfway, round to even *)
    IF ABS(d - r) = 0.5d0 AND
       FLOAT(ROUND(r / 2.0d0), LONGREAL) * 2.0d0 # r THEN
      RETURN FromLR(r - FLOAT(ORD(d > 0.0d0) * 2 - 1, LONGREAL))
    ELSE
      RETURN FromLR(r)
    END
  END
\end{lstlisting}

\noindent Alternatively, use {\tt FLOOR(d + 0.5)} with an even-check
when {\tt d} is a half-integer.

\subsection{{\tt remainder} Sign Convention}

R5RS specifies that {\tt remainder} returns a result with the sign
of the \emph{dividend} (first argument).  MScheme used Modula-3's
{\tt MOD}, which gives the \emph{divisor}'s sign.
\textbf{Already fixed} (see {\tt SchemePrimitive.m3}).

%======================================================================
\section{Summary of File Changes}
\label{sec:impl-summary}
%======================================================================

Table~\ref{tab:files} lists all files created or modified in Phase~1.

\begin{table}[htbp]
\centering
\begin{tabular}{llp{6cm}}
\toprule
\textbf{File} & \textbf{Action} & \textbf{Change} \\
\midrule
{\tt SchemeInt.i3} & New & Exact integer interface \\
{\tt SchemeInt.m3} & New & Fixnum cache, overflow-checked
  arithmetic, bignum promotion \\
{\tt SchemeLongReal.m3} & Modify & {\tt FromO}, {\tt Int}:
  add {\tt REF INTEGER} and {\tt Mpz.T} branches \\
{\tt SchemeInputPort.m3} & Modify & Reader: integer vs.\ float
  dispatch; {\tt +inf.0}/{\tt +nan.0} syntax \\
{\tt SchemeUtils.m3} & Modify & Printer: {\tt REF INTEGER} and
  {\tt Mpz.T} branches; {\tt eqv?}\ exactness;
  R7RS Inf/NaN output \\
{\tt SchemePrimitive.m3} & Modify & Predicates, arithmetic
  dispatch, comparison dispatch, banker's rounding \\
{\tt compiler.scm} & Modify & {\tt constant-ref}, inline
  arithmetic, inline predicates \\
{\tt mscheme/src/m3makefile} & Modify & Add {\tt Module("SchemeInt")},
  {\tt import("mpz")} \\
\bottomrule
\end{tabular}
\caption{Files created or modified in Phase~1.}
\label{tab:files}
\end{table}

%======================================================================
\section{Phase~1 Test Suite}
\label{sec:impl-tests}
%======================================================================

A test-first approach: the file {\tt mscheme-doc/test-numeric-tower.scm}
contains 347~tests that define the target behavior for Phase~1.
Run with:
\begin{lstlisting}
mscheme test-numeric-tower.scm
\end{lstlisting}

\noindent Against the current (pre-tower) MScheme, 287 tests pass
and 60~fail.  The failures fall into these categories, each
corresponding to a subsystem that Phase~1 must modify:

\begin{table}[htbp]
\centering
\begin{tabular}{lcp{6.5cm}}
\toprule
\textbf{Category} & \textbf{Fails} & \textbf{Root Cause} \\
\midrule
{\tt inexact?}\ predicates &
  10 & Integer-valued floats ({\tt 0.0}, {\tt 1.0}, {\tt 1e10})
       not recognized as inexact \\
{\tt eqv?}\ across exactness &
  2 & {\tt (eqv?\ 1 1.0)} returns {\tt \#t}; should be {\tt \#f} \\
Exactness contagion &
  6 & {\tt (+ 1 2.0)} not recognized as inexact \\
Bignum exactness &
  15 & {\tt (expt 2 62)} et al.\ return {\tt Mpz.T} but
      {\tt exact?}\ does not recognize it \\
Fixnum cache {\tt eq?} &
  4 & No fixnum cache; {\tt REF INTEGER} freshly allocated \\
{\tt \#i} prefix &
  1 & {\tt \#i1} does not force inexact \\
Bignum precision &
  5 & Bignums compared/printed via {\tt LONGREAL} lose
      precision (e.g.\ {\tt 2\^{}100 = 2\^{}100+1}) \\
Banker's rounding &
  4 & {\tt ROUND} uses half-away-from-zero, not half-to-even \\
Division by zero &
  4 & Returns {\tt Inf} instead of raising an error \\
Bignum {\tt gcd}/{\tt lcm} &
  4 & Overflow in {\tt TRUNC} when operands exceed
      {\tt LAST(INTEGER)} \\
30! exactness &
  3 & Factorial computed via {\tt LONGREAL}; loses precision
      beyond $2^{53}$ \\
Bignum print round-trip &
  2 & Bignums print in exponential notation via {\tt LONGREAL} \\
\bottomrule
\end{tabular}
\caption{Test failure breakdown against pre-tower MScheme (287/347 passing).}
\label{tab:failures}
\end{table}

\noindent The tests are organized in 24~sections:

\begin{enumerate}[nosep]
\item Exact/inexact predicates
\item Number type predicates ({\tt number?}, {\tt integer?}, {\tt real?}, etc.)
\item Printing (exact integers without decimal point)
\item {\tt eqv?}\ exactness semantics
\item Arithmetic exactness preservation and contagion
\item Exact integer arithmetic (basic operations)
\item Division semantics (exact vs.\ inexact result)
\item {\tt gcd} and {\tt lcm}
\item Comparisons across exactness
\item {\tt exact->inexact} / {\tt inexact->exact}
\item Bit-boundary rollover (31/32/63/64-bit boundaries, {\tt LONGREAL} precision)
\item Infinity and NaN (IEEE~754 special values)
\item Rounding procedures ({\tt floor}, {\tt ceiling}, {\tt truncate}, {\tt round})
\item Fixnum cache identity ({\tt eq?})
\item Reader exactness prefixes ({\tt \#e}, {\tt \#i})
\item Non-decimal base literals ({\tt \#b}, {\tt \#o}, {\tt \#x})
\item {\tt number->string} / {\tt string->number}
\item Edge cases (zero, {\tt max}/{\tt min}, {\tt abs})
\item Error cases (type errors, division by zero)
\item Negative bignum arithmetic
\item Mixed fixnum/bignum operations
\item Bignum comparisons
\item Bignum print/read round-trip
\item {\tt expt} edge cases ({\tt 0\^{}0}, {\tt -1\^{}$n$}, negative exponents)
\item {\tt min}/{\tt max} with bignums
\item Zero/one-arg arithmetic (negation, reciprocal)
\end{enumerate}

\noindent As each subsystem is modified, the corresponding test
sections should go from failing to passing.  When all 347~tests
pass, Phase~1 is complete.

\end{document}
